\subsection*{Категориальные признаки в деревьях. Как быть?}

Значит, есть 3 варианта как быть с категориальными признаками:

\begin{itemize}
  \item Для категориальных признаков строим для дерева не 2 потомка, а столько, сколько значений может быть у категориального признака.
\end{itemize}

К слову, такой подход называется \textit{multi-way split};

\begin{itemize}
  \item Пусть \(Q = u_1, \ldots, u_q\) — множество значений категориального множества. Делаем разбиение на 2 непересекающихся множества: \(Q = Q_1 \sqcup Q_2\).
\end{itemize}

а предикат устанавливаем в виде \(b_j = [x_j \in Q]\). Однако это очень долго (придётся перебирать \(\frac{2^q- 2}{2}\) вариантов);

\begin{itemize}
  \item А третий вариант стоит рассмотреть как отдельную часть. Для примера рассмотрим задачу бин. классификации, упорядочим все значения категориального признака по следующему принципу:
\end{itemize}

\[
\frac{1}{N_{u_i}} \sum_{x_j \in R(u_i))} [y_i = +1]
\]
Где \(R(u_i)\) — те объекты из вершины \(R\), в которых значение категориального признака равно \(u_i\), а \((N_{u_i})\) — их количество. 

Т.е. мы просто считаем долю объектов, которые имеют конкретное значение рассматриваемого категориального признака и при этом имеют положительный класс. Получившуюся последовательность сортируем по возрастанию.

Далее применим простое отображение: заменим первый элемент получившейся последовательности на 1, второй на 2 и так далее, т.е. \(u_{i(1)} \to 1, u_{i(2)} \to 2, \ldots, u_{i(q)} \to q\). Всё что нам осталось сделать — работать с этим в будущем как с числовым признаком.
